% =====================================================================
%  ESERCITAZIONE DIAGRAMMI DI FLUSSO — Sistemi Automatici 3BETC
%  In preparazione alla verifica del 24/10/2026 (esercizio 4: ciclo).
%  Imposta questo file come Main document e compila: un solo PDF con
%     pag. 1 Esercizio 1 (più semplice)   pag. 2 Esercizio 2 (più impegnativo)
%     pag. 3 Soluzione dell'esercizio 1   pag. 4 Soluzione dell'esercizio 2
%     ogni esercizio ha una tabella di traccia da completare (prova con valori dati)
% =====================================================================
\documentclass[a4paper,11pt]{article}
\usepackage{stili/comune}
\usepackage[scaled=0.92]{helvet}
\renewcommand{\familydefault}{\sfdefault}
\usepackage[a4paper,top=2.3cm,bottom=2.2cm,left=2cm,right=2.2cm]{geometry}
\pagestyle{plain}
\setlength{\parindent}{0pt}
\setlength{\parskip}{5pt}

\newcommand{\titolo}[1]{%
  \begin{center}
    {\small Sistemi Automatici -- 3BETC -- Esercitazione sui diagrammi di flusso\par}
    \vspace{4pt}{\Large\bfseries #1\par}
  \end{center}\vspace{6pt}}

\begin{document}

% ======================= TESTI ========================================
\titolo{Esercitazione sui diagrammi di flusso}

\textbf{Esercizio 1}

In una cisterna ci sono dei litri d'acqua. Con quest'acqua si riempiono bottiglie tutte uguali,
una alla volta. Disegna il diagramma di flusso di un algoritmo che conta quante bottiglie si
possono riempire completamente.

\textbf{Dati da leggere:} $Q$ = litri d'acqua nella cisterna; $B$ = litri contenuti in una bottiglia.\\
\textbf{Altre variabili:} $N$ = numero di bottiglie riempite.

Si riempiono bottiglie finché nella cisterna c'è abbastanza acqua per riempirne una intera.
Alla fine visualizza $N$ e i litri rimasti $Q$.

\textbf{Prova:} con $Q = 17$ e $B = 4$ segui il diagramma e completa la tabella con i valori dopo
ogni giro del ciclo, fino a quando il ciclo finisce.
\begin{center}\renewcommand{\arraystretch}{1.35}
\begin{tabular}{|c|p{4.2cm}|p{2.6cm}|p{2.6cm}|}\hline
  Giro & Condizione: vera o falsa? & $Q$ & $N$\\\hline
  Inizio & & &\\\hline
  1 & & &\\\hline
  2 & & &\\\hline
  3 & & &\\\hline
  4 & & &\\\hline
  5 & & &\\\hline
  6 & & &\\\hline
\end{tabular}\end{center}

\newpage
\textbf{Esercizio 2}

Un serbatoio, inizialmente vuoto, si riempie versando un secchio d'acqua alla volta. Il serbatoio
però ha una piccola perdita: dopo ogni secchio versato escono $D$ litri. Disegna il diagramma di
flusso di un algoritmo che conta quanti secchi servono per riempirlo.

\textbf{Dati da leggere:} $C$ = capacità del serbatoio (in litri); $S$ = litri contenuti in un
secchio; $D$ = litri persi dopo ogni secchio.\\
\textbf{Altre variabili:} $L$ = litri presenti nel serbatoio; $N$ = numero di secchi versati.

Si versano secchi finché il serbatoio non è pieno. Se però sono già stati versati 50 secchi e il
serbatoio non è ancora pieno, l'algoritmo si ferma e visualizza ``Perdita troppo grande''.
Se il serbatoio si riempie, alla fine visualizza $N$ e $L$.

\textbf{Prova:} con $C = 100$, $S = 30$ e $D = 5$ segui il diagramma e completa la tabella con i
valori dopo ogni giro del ciclo, fino a quando il ciclo finisce. Che cosa visualizza l'algoritmo?
\begin{center}\renewcommand{\arraystretch}{1.35}
\begin{tabular}{|c|p{3.2cm}|p{3.2cm}|p{2.2cm}|p{2.2cm}|}\hline
  Giro & Prima condizione & Seconda condizione & $L$ & $N$\\\hline
  Inizio & & & &\\\hline
  1 & & & &\\\hline
  2 & & & &\\\hline
  3 & & & &\\\hline
  4 & & & &\\\hline
  5 & & & &\\\hline
  6 & & & &\\\hline
\end{tabular}\end{center}

\newpage
% ======================= SOLUZIONE 1 ==================================
\titolo{Soluzione dell'esercizio 1}

\begin{center}
\begin{tikzpicture}[x=1cm,y=1cm]
  \node[fcterm] (ini)  at (0,0) {Inizio};
  \node[fcio]   (in)   at (0,-1.1) {Leggi $Q$, $B$};
  \node[fcinit] (ini0) at (0,-2.2) {$N \gets 0$};
  \node[fcdec]  (d1)   at (0,-3.75) {$Q \geq B$};
  \node[fcproc] (agg)  at (0,-5.3) {$Q \gets Q - B$};
  \node[fcproc] (inc)  at (0,-6.4) {$N \gets N + 1$};
  \node[fcio]   (out)  at (5.0,-3.75) {Scrivi $N$, $Q$};
  \node[fcterm] (fin)  at (5.0,-4.95) {Fine};
  \draw[fcfreccia] (ini) -- (in);
  \draw[fcfreccia] (in) -- (ini0);
  \draw[fcfreccia] (ini0) -- (d1);
  \draw[fcfreccia] (d1) -- node[fcsi,right] {Sì} (agg);
  \draw[fcfreccia] (agg) -- (inc);
  \draw[fcfreccia] (inc.south) -- ++(0,-0.5) -| (-2.4,-2.95) -- (0,-2.95);
  \draw[fcfreccia] (d1) -- node[fcno,above] {No} (out);
  \draw[fcfreccia] (out) -- (fin);
\end{tikzpicture}
\end{center}

\textbf{Come funziona.}
\begin{itemize}[nosep]
  \item $N \gets 0$ è l'inizializzazione (esagono): va fatta una volta sola, \emph{prima} del ciclo.
  \item Il rombo $Q \geq B$ è la condizione del ciclo: se è vera si riempie un'altra bottiglia,
        se è falsa il ciclo finisce.
  \item Ogni giro del ciclo è una bottiglia riempita: $Q$ cala di $B$, $N$ aumenta di 1, poi si
        torna al rombo.
  \item Si usa $\geq$ e non $>$: se nella cisterna restano esattamente $B$ litri, si riempie
        ancora una bottiglia.
\end{itemize}

\textbf{Prova} ($Q = 17$, $B = 4$):

\begin{center}
\renewcommand{\arraystretch}{1.15}
\begin{tabular}{|c|c|c|c|}\hline
  Giro & Condizione $Q \geq B$ & $Q$ & $N$\\\hline
  Inizio & --  & 17 & 0\\\hline
  1   & $17 \geq 4$: vera & 13 & 1\\\hline
  2   & $13 \geq 4$: vera & 9  & 2\\\hline
  3   & $9 \geq 4$: vera  & 5  & 3\\\hline
  4   & $5 \geq 4$: vera  & 1  & 4\\\hline
  --  & $1 \geq 4$: falsa & \multicolumn{2}{c|}{fine del ciclo}\\\hline
\end{tabular}
\end{center}
Si visualizza $N = 4$ bottiglie e $Q = 1$ litro rimasto.

\newpage


% ======================= SOLUZIONE 2 ==================================
\titolo{Soluzione dell'esercizio 2}

\begin{center}
\begin{tikzpicture}[x=1cm,y=1cm]
  \node[fcterm] (ini)  at (0,0) {Inizio};
  \node[fcio]   (in)   at (0,-1.1) {Leggi $C$, $S$, $D$};
  \node[fcinit] (ini0) at (0,-2.2) {$L \gets 0$, \ $N \gets 0$};
  \node[fcdec]  (d1)   at (0,-3.75) {$L < C$};
  \node[fcdec]  (d2)   at (0,-5.45) {$N < 50$};
  \node[fcproc] (agg)  at (0,-7.0) {$L \gets L + S - D$};
  \node[fcproc] (inc)  at (0,-8.1) {$N \gets N + 1$};
  \node[fcio]   (out1) at (4.9,-3.75) {Scrivi $N$, $L$};
  \node[fcio]   (out2) at (4.9,-5.45) {Scrivi ``Perdita troppo grande''};
  \node[fcterm] (fin)  at (4.9,-7.0) {Fine};
  \draw[fcfreccia] (ini) -- (in);
  \draw[fcfreccia] (in) -- (ini0);
  \draw[fcfreccia] (ini0) -- (d1);
  \draw[fcfreccia] (d1) -- node[fcsi,right] {Sì} (d2);
  \draw[fcfreccia] (d2) -- node[fcsi,right] {Sì} (agg);
  \draw[fcfreccia] (agg) -- (inc);
  \draw[fcfreccia] (inc.south) -- ++(0,-0.5) -| (-2.5,-2.95) -- (0,-2.95);
  \draw[fcfreccia] (d1) -- node[fcno,above] {No} (out1);
  \draw[fcfreccia] (d2) -- node[fcno,above] {No} (out2);
  \draw[fcfreccia] (out2) -- (fin);
  \draw[fcfreccia] (out1.east) -- (7.9,-3.75) |- (fin.east);
\end{tikzpicture}
\end{center}

\textbf{Come funziona.}
\begin{itemize}[nosep]
  \item Nel ciclo ci sono \emph{due} controlli: prima si guarda se il
        serbatoio è pieno ($L < C$?), poi se si è arrivati al limite di 50 secchi ($N < 50$?).
        Si versa un altro secchio solo se tutte e due le risposte sono Sì.
  \item In ogni giro il livello aumenta di $S$ (il secchio) e diminuisce di $D$ (la perdita):
        $L \gets L + S - D$.
  \item Le due uscite dal ciclo portano a messaggi diversi, ma finiscono entrambe in Fine.
  \item Il limite di 50 secchi serve perché, se $D \geq S$, il serbatoio non si riempirebbe mai e
        il ciclo non finirebbe più.
\end{itemize}

\textbf{Prova} ($C = 100$, $S = 30$, $D = 5$): a ogni secchio il livello sale di $30 - 5 = 25$ litri.

\begin{center}
\renewcommand{\arraystretch}{1.15}
\begin{tabular}{|c|c|c|c|c|}\hline
  Giro & Condizione $L < C$ & Condizione $N < 50$ & $L$ & $N$\\\hline
  Inizio & --  & --  & 0  & 0\\\hline
  1   & $0 < 100$: vera  & $0 < 50$: vera & 25  & 1\\\hline
  2   & $25 < 100$: vera & $1 < 50$: vera & 50  & 2\\\hline
  3   & $50 < 100$: vera & $2 < 50$: vera & 75  & 3\\\hline
  4   & $75 < 100$: vera & $3 < 50$: vera & 100 & 4\\\hline
  --  & $100 < 100$: falsa & & \multicolumn{2}{c|}{fine del ciclo}\\\hline
\end{tabular}
\end{center}
Si visualizza $N = 4$ secchi e $L = 100$ litri (il serbatoio è pieno esattamente). Con $S = 5$ e
$D = 5$, invece, il livello resta sempre 0: dopo 50 secchi si visualizza ``Perdita troppo grande''.

\end{document}
